\documentclass[11pt]{article}
\usepackage[margin=1in]{geometry}
\usepackage{amsmath,amssymb,amsthm}
\newtheorem{theorem}{Theorem}
\title{A disproof of Conjecture 00000000748}
\author{earthking11}
\date{October 6, 2026}
\begin{document}
\maketitle
\section*{Statement and counterexample}
The conjecture classifies the polynomials compatible with the Teichmuller lift
$w:\mathbb F_p\to\mathbb Z_p$, claiming that the only minimum-degree examples
satisfying $P\circ w=w\circ\overline P$ are identity and Frobenius.  As written,
constant polynomials are not excluded.

\begin{theorem}
The stated classification is false.
\end{theorem}
\begin{proof}
Take $P(X)=0\in\mathbb Z_p[X]$ and $\overline P(X)=0\in\mathbb F_p[X]$.
The defining property $w(0)=0$ gives, for every $x\in\mathbb F_p$,
\[
 P(w(x))=0=w(0)=w(\overline P(x)).
\]
Hence the required commuting identity holds.  The polynomial has minimum
possible degree (degree zero among nonempty coefficient lists; under the
standard convention the zero polynomial has degree $-\infty$).  Nevertheless
it is neither identity nor Frobenius: at $1$ both listed maps have value $1$,
whereas $P(1)=0$.
\end{proof}

The same argument works with the constant-one polynomial.  Therefore any valid
classification must at least add a nonconstant or normalization hypothesis;
neither occurs in the filed conjecture.

\section*{Lean 4 certificate}
The project in \texttt{lean4/} defines the commuting-square predicate
\texttt{LiftCompatible}.  The theorem \texttt{counterexample} constructs the
zero maps, proves compatibility from $w(0)=0$, and proves by evaluation at $1$
that the map is neither identity nor Frobenius.  It contains no
\texttt{sorry}, \texttt{native\_decide}, or additional axiom.
\end{document}
